\documentclass[letterpaper,10pt]{amsart}
\usepackage{graphicx}
\usepackage{latexsym}
\usepackage{amsmath}
\usepackage{amsfonts}
\usepackage{amsthm}
\usepackage{amssymb}
\usepackage{tikz-cd}
\usepackage[T1]{fontenc}
\usepackage{lmodern,microtype}
\usepackage{needspace}
\usepackage[hidelinks]{hyperref}
\hypersetup{pdftitle={Ranges and inseparability in recursion categories -- Version 2},
pdfauthor={Florian Lengyel}}

\theoremstyle{plain}
\newtheorem{thm}{Theorem}[section]
\newtheorem{cor}[thm]{Corollary}
\newtheorem{lem}[thm]{Lemma}
\newtheorem{claim}[thm]{Claim}
\newtheorem{axiom}[thm]{Axiom}
\newtheorem{conj}[thm]{Conjecture}
\newtheorem{fact}[thm]{Fact}
\newtheorem{hypo}[thm]{Hypothesis}
\newtheorem{assum}[thm]{Assumption}
\newtheorem{prop}[thm]{Proposition}
\newtheorem{crit}[thm]{Criterion}
\theoremstyle{definition}
\newtheorem{defn}[thm]{Definition}
\newtheorem{exmp}[thm]{Example}
\newtheorem{rem}[thm]{Remark}
\newtheorem{prob}[thm]{Problem}
\newtheorem{prin}[thm]{Principle}
\newtheorem{alg}[thm]{Algorithm}
\newcommand{\map}[2]{{#1}\rightarrow{#2}}
\newcommand{\MAP}[3]{{#1}:{#2}\rightarrow{#3}}
\newcommand{\sento}[2]{{#1}\mapsto{#2}}
\newcommand{\rst}[1]{\overline{#1}}

\newcommand{\XP}{X^\#}
\newcommand{\IBAR}{\overline{i}_0}
\newcommand{\TUR}[1]{\mathrm{Tur}\left({#1}\right)}
\newcommand{\IT}[1]{\mathrm{It}\left({#1}\right)}

\begin{document}
\title{Ranges and inseparability in recursion categories}
\author{Florian Lengyel}
\date{Version 2\\1 November 2026}
\thanks{Email address: flengyel@gc.cuny.edu}
\address{The Graduate School and University Center, %
Ph.D Program in Computer Science, %
The City University of New York, 365 Fifth Avenue, New York, NY %
10016}

\subjclass[2020]{03D75, 03G30}

\begin{abstract}
In a $P$-recursion category with binary coproducts and ranges, the
domain of a weak Turing morphism applied to the diagonal has no
complement when the identity of its underlying object is nonzero.
We also construct recursively inseparable domains using a range functor
to a split restriction category satisfying weak choice, in which inverse
images of domains along every morphism preserve binary joins.
This construction requires each target morphism selecting between the
images of a complemented source domain and its complement to have a preimage under
the functor. Joins of source morphisms with disjoint domains suffice
for this requirement. A counterexample satisfies the remaining
hypotheses, but the constructed pair is separable and the required
preimage does not exist.

\end{abstract}

\maketitle

\section{Introduction}
Di Paola and Heller developed recursion categories as an algebraic
setting for recursion theory, with the aim of obtaining a categorical
formulation of incompleteness \cite{MR88i:03078}. Morphisms represent
partial computations, and restriction idempotents represent their
domains. Rosolini's $P$-categories provide a more general setting for
these constructions \cite{coneit,MR90e:18001}, used in subsequent
existence and relativization theorems
\cite{MR92e:03067,MR95c:03110,leng}.

We study two diagonal arguments in this setting. The first produces a
domain with no complement in a $P$-recursion category with binary
coproducts and ranges, provided that the identity of the underlying
object is nonzero (Theorem~\ref{ThmNoncomp}). The second produces a
recursively inseparable pair of domains. Such a pair consists of two
disjoint domains for which no complemented domain contains the first
while its complement contains the second.

The inseparability argument uses a range functor, which preserves the
domain and range operations, to a split restriction category satisfying
weak choice, in which inverse images of domains along every morphism
preserve binary joins. Under the target hypotheses of
Theorem~\ref{ThmRecIns}, every complemented source domain determines a
target morphism that agrees with one coproduct injection on the image
of the domain and with the other on the image of its complement.
The lifting hypothesis requires that this morphism
be the image of a source morphism. Joins of source morphisms with
disjoint domains suffice for this hypothesis
(Proposition~\ref{PropDisjointJoins}). The counterexample in
Section~\ref{sec:counterexample} satisfies the other hypotheses of the
theorem, but its diagonal pair is separable; the target morphism has
no source preimage.

\Needspace{6\baselineskip}
Examples defined by arithmetic provability motivate the use of a
comparison functor. Let $T$ be a consistent
recursively enumerable extension of Peano arithmetic ($\mathrm{PA}$).
In the category $S_T$, morphisms are equivalence classes of program
indices: two indices are equivalent precisely when $T$ proves a single
universally quantified statement asserting that the programs have the
same domain and values \cite[p.~646]{MR93b:03103}. Montagna's category
$S$ is the case $T=\mathrm{PA}$ \cite[p.~105]{MR90i:03046}.
The proposed system $S'_T$, with $S'=S'_{\mathrm{PA}}$, instead uses
separate proofs of equality at each standard input. Pointwise provable
equality is not preserved by composition, so these proposed quotients
do not define categories under that composition
\cite[Corollary~2]{leng2026composition}.

Di Paola and Montagna prove that $S_T$ is not locally connected for
any such $T$ \cite[Corollary~4.4]{MR93b:03103}; if $T$ is not
$\Sigma^0_1$-sound, it has no nonzero connected domains
\cite[Theorem~4.3(2)]{MR93b:03103}. Heller's existence theorem uses
local connectedness \cite{MR92e:03067}. These examples therefore
motivate extending that theorem beyond locally connected categories,
as proposed in \cite{MR93b:03103} and developed using range functors
in \cite{leng}.

Section~\ref{sec:preliminaries} introduces restriction categories,
near products, and Turing morphisms. Section~\ref{sec:diagonal}
develops the range calculus needed for the noncomplementation theorem
and proves that theorem. Section~\ref{sec:inseparability} constructs
the target morphisms, states the lifting hypothesis, and proves the
inseparability theorem, its sufficient condition, and the counterexample.

\section{Restriction categories and recursion categories}\label{sec:preliminaries}

\subsection{Restriction categories}
We write $gf$ for the composite that applies $f$ first and then $g$.
A morphism $e:X\to X$ is \textit{idempotent} if $e^2=e$.

A \textit{restriction category} is a category $\mathbf C$ with an
operation assigning to each morphism $f:X\to Y$ an endomorphism
$\rst f:X\to X$, subject to the following axioms whenever the
composites are defined \cite{MR1871071,MR1963657}.

\noindent\textbf{R1.} $f\rst f=f$.

\noindent\textbf{R2.} $\rst f\rst g=\rst g\rst f$.

\noindent\textbf{R3.} $\rst{g\rst f}=\rst g\rst f$.

\noindent\textbf{R4.} $\rst g f=f\rst{gf}$.

These axioms imply that $\rst f$ is idempotent and that
$\rst{gf}=\rst{\rst g f}$. The endomorphisms $\rst f$ are called
\textit{restriction idempotents}, or \textit{domains}. For an object
$X$, write $\mathrm{Dom}(X)$ for its domains. They form a
meet-semilattice: $\alpha\le\beta$ means $\alpha\beta=\alpha$, and
$\alpha\wedge\beta=\alpha\beta$. A morphism $f:X\to Y$ is
\textit{total} if $\rst f=1_X$; total morphisms form a subcategory
$\mathbf C_T$.

For parallel morphisms $f,g:X\to Y$, the \textit{restriction order}
is defined by $f\le g$ if $f=g\rst f$. It restricts to the order
just defined on domains. A join $f\vee g$, when it exists, is the
least upper bound in this order.

A \textit{system of zero morphisms} consists of morphisms
$0_{X,Y}:X\to Y$ such that $g0_{X,Y}f=0_{W,Z}$ for all
$f:W\to X$ and $g:Y\to Z$. Such a system is unique when it exists.
We write $0_X=0_{X,X}$ and use \textit{restriction zeros}, which
satisfy $\rst{0_{X,Y}}=0_X$. Domains $\alpha,\beta$ on $X$ are
\textit{disjoint} if $\alpha\beta=0_X$; parallel morphisms $f,g$
have \textit{disjoint domains} if $\rst f\rst g=0_X$.
A morphism $f:X\to Y$ is \textit{weakly total} if $fg=0$ implies
$g=0$ for every $g:Z\to X$.

A restriction category is \textit{split} if each domain
$\varepsilon:X\to X$ has a splitting $\varepsilon=me$, where
$e:X\to Z$, $m:Z\to X$, and $em=1_Z$. The map $m$ is then a
total monomorphism.

Binary \textit{restriction coproducts} are categorical coproducts with
total injections $i_0:X\to X+Y$ and $i_1:Y\to X+Y$. For morphisms
$f:X\to Z$ and $g:Y\to Z$, write $[f,g]:X+Y\to Z$ for their
copairing. The \textit{codiagonal} is
$\nabla_X=[1_X,1_X]:X+X\to X$; thus $[f,g]=\nabla_Z(f+g)$.
Restrictions satisfy $\rst{f+g}=\rst f+\rst g$.

Cockett and Lack's restriction categories \cite{MR1871071,MR1963657},
equivalently Stefani's RDP-categories \cite{MR2003a:18003}, describe
domains by this unary operation. We next recall the near products and
diagonals needed to define Turing morphisms.

\subsection{\texorpdfstring{$P$}{P}-categories}
Let $\mathbf{C}$ be a category. The {\it diagonal functor}
$\MAP{\Delta_\mathbf{C}}{\mathbf C}{\mathbf{C}\times\mathbf{C}}$
is given on objects by $\sento{A}{(A,A)}$ and on morphisms by
$\sento{f}{(f,f)}$.

A {\it P-category} consists of a category $\mathbf{C}$ together
with a bifunctor $\MAP{\times}{\mathbf{C}\times\mathbf{C}}{\mathbf
C}$, called a {\it near product}, a natural transformation
$\MAP{\Delta}{1_\mathbf{C}}{\times\circ\Delta_{\mathbf C}}$, and
for each object $X$ of $\mathbf{C}$, natural transformations
\[
\begin{array}{ll}
p_{0 (\ ),X}:{(\ )\times X}\rightarrow{1_{\mathbf{C}}}, & 
p_{1 X, (\ )}:{X\times (\ ) }\rightarrow{1_{\mathbf{C}}}
\end{array}
\]
such that the conditions i), ii) and iii) below are satisfied ~\cite{coneit}.

For morphisms $f:X\to Y$ and $g:X\to Z$, put
$\langle f,g\rangle=(f\times g)\Delta_X:X\to Y\times Z$.

i) The following equations hold.
\[
\begin{array}{ll} %
p_{0 X,X}\Delta_X=1_X=p_{1 X, X}\Delta_X, & (p_{0 X,Y}\times
p_{1 X,Y})\Delta_{X\times Y}=1_{X\times Y},\\
p_{0 X,Y}(1_X\times p_{0 Y,Z})=p_{0 X,Y\times Z}, & p_{0
X,Z}(1_X\times p_{1 Y,Z})=p_{0 X,Y\times Z},\\
p_{1 X,Z}(p_{0 X,Y}\times 1_Z)=p_{1 X\times Y, Z}, & p_{1
Y,Z}(p_{1 X,Y}\times1_Z )=p_{1 X\times Y, Z}.
\end{array}
\]

ii) There is an isomorphism of functors
$\mathrm{ass}_\times:((-\times-)\times-)\Rightarrow(-\times(-\times-))$
from $\mathbf C^3$ to $\mathbf C$, whose component
$\mathrm{\mathrm{ass}_\times}_{X,Y,Z}$  is given by
\begin{align*}
\MAP{{\langle p_{0 X,Y} p_{0 X\times Y,Z}, %
      \langle p_{1 X,Y} p_{0 X\times Y,Z}, %
       p_{1 X\times Y,Z}\rangle\rangle }}{ (X \times Y )\times Z}{ X\times (Y \times Z)},%
\end{align*}
iii) Let $T$ be the endofunctor on
$\mathbf{C}\times\mathbf{C}$ given by $\sento{(X,Y)}{(Y,X)}$.
There is an isomorphism
$\mathrm{tr}_\times:\times\Rightarrow\times\circ T$ of functors
from $\mathbf C^2$ to $\mathbf C$, with component
$\mathrm{tr}_{\times X,Y}=\langle p_{1X,Y},p_{0X,Y}\rangle:
X\times Y\to Y\times X$.
The isomorphisms $\mathrm{\mathrm{ass}_\times}$ and $\mathrm{tr}_\times$ are 
natural in each of their component variables. 

The induced restriction of a morphism $\varphi:X\to Y$ is
$\rst\varphi=p_0(1_X\times\varphi)\Delta_X
=p_1(\varphi\times1_X)\Delta_X$. This operation satisfies
R1--R4 \cite{MR88i:03078,coneit}.

The near product restricts to a categorical product on
$\mathbf C_T$ \cite{coneit}. 

In general, a $P$-category is equivalent to one obtained from a
category with products, following a procedure analogous to that by
which the category $\mathbf{Pfn}$ of sets and partial functions
can be obtained from the category $\mathbf{Set}$ of sets 
~\cite[Theorem 1.2]{MR90e:18001}.

The projection $p_0$ is natural in the first variable, and $p_1$ is
natural in the second. Thus $p_0(\varphi\times1_X)=\varphi p_0$,
whereas $p_1(\varphi\times1_X)=1_Xp_1$ can fail: the right-hand
side is total, while the left-hand side need not be. In
general, the projections $p_{i,X,Y}$ of a $P$-category
$\mathbf{C}$ are natural in both variables if and only if for any morphisms
$\MAP{\varphi_i}{Y}{X_i}$ ($i=0,1$) in 
$\mathbf{C}$, $p_i\langle \varphi_0,\varphi_1\rangle = \varphi_i$.
Cockett and Lack observe that $P$-categories are ``symmetric monoidal categories
in which each object has a monoidal natural cocommutative coassociative
comultiplication (and possibly an unnatural counit)'' ~\cite{MR1871071}.

A $P$-category \textit{has coproducts} if its induced restriction
category has binary restriction coproducts. It is \textit{prodominical}
if it has zero morphisms and $\phi\times0_{C,D}=0_{A\times C,B\times D}$
for every morphism $\phi:A\to B$ \cite{MR92e:03067}. We assume
throughout that our $P$-categories are prodominical. Their zero domains
are the least elements of $\mathrm{Dom}(X)$.

A prodominical category is \textit{dominical} if every weakly total
morphism is total and $f\times g=0$ implies $f=0$ or $g=0$
\cite[p.~36]{leng}.

A category is \textit{isotypical} if any two of its objects are isomorphic.

\subsection{Turing morphisms and \texorpdfstring{$P$}{P}-recursion categories}
A {\it Turing morphism} in a prodominical isotypical $P$-category is a morphism %
$\MAP{\tau}{W\times X}{Y}$ such that for any $\MAP{\phi}{V\times
X}{Y}$ there exists a total $\MAP{g}{V}{W}$, such that the
following diagram commutes.
\[
\begin{tikzcd}
V\times X \arrow[r,"g\times1_X"] \arrow[dr,"\phi"']
  & W\times X \arrow[d,"\tau"] \\
  & Y
\end{tikzcd}
\]

A {\it $P$-recursion category} (or, simply, a {\it recursion
category}) is a prodominical isotypical $P$-category  which contains a Turing
morphism. Previously, the term {\it recursion category}
was used for dominical $P$-recursion categories; for us a recursion category
is always a $P$-recursion category. The prototypical example of a Turing morphism $\MAP{t}{X\times X}{X}$ is the universal partial recursive function ~\cite[Theorem II.1.8, page 132]{MR90d:03072}. Examples of Turing morphisms and recursion categories are given in ~\cite{MR88i:03078,MR90i:03046,MR92e:03067,MR93b:03103,MR95c:03110,leng}.

A morphism $t:X\times X\to X$ is \textit{weak Turing} if, for every
$f:X\to X$, there is a total morphism $e:X\to X$ such that
\begin{equation}\label{DefWeak}
\begin{tikzcd}
X\times X \arrow[r,"e\times1_X"] \arrow[d,"p_{1,X,X}"']
  & X\times X \arrow[d,"t"] \\
X \arrow[r,"f"'] & X
\end{tikzcd}
\end{equation}
commutes. Such an $e$ is a \textit{$t$-index} of $f$. The morphism $t$
is \textit{idempotent weak Turing} if every $f$ has an idempotent
$t$-index. Every Turing morphism of type $X\times X\to X$ is weak
Turing; conversely, a Turing morphism can be obtained from a weak Turing
morphism \cite[Lemma 9.1]{MR92e:03067}.

For every total $t$-index $e$ of $f$, precomposing
\eqref{DefWeak} with $\langle1_X,e\rangle$ gives
\begin{equation}\label{EqDiagonalIndex}
t\Delta_Xe=fe.
\end{equation}
Indeed, $(e\times1_X)\langle1_X,e\rangle=\Delta_Xe$ and
$p_{1,X,X}\langle1_X,e\rangle=e$. This identity will be used in
both diagonal arguments.

\begin{prop}\label{IdempotentTheorem}
Every $P$-recursion category $\mathbf C$ has an object $Y$, an
idempotent weak Turing morphism $t:Y\times Y\to Y$, and a function
$N:\mathbf C_T(Y,Y)\to\mathbf C_T(Y,Y)$ such that, for every
morphism $\varphi:Y\to Y$ and every total $t$-index $e:Y\to Y$ of
$\varphi$, the morphism $N(e)$ is an idempotent $t$-index of $\varphi$.
\end{prop}
\begin{proof}
Choose a Turing morphism $\tau:X\times X\to X$, put $Y=X\times X$,
and choose an isomorphism $a:Y\to X$. Let $q_0,q_1:Y\to X$ be
the projections. The morphism
$t=a^{-1}\tau(q_1\times a):Y\times Y\to Y$ is Turing. Given an
object $V$ and a morphism $f:V\times Y\to Y$, choose a total
morphism $c:V\to X$ such that
\[
 \tau(c\times1_X)=af(1_V\times a^{-1}).
\]
Then $g=\Delta_Xc:V\to Y$ is total, $q_1g=c$, and
\[
 t(g\times1_Y)=a^{-1}\tau(c\times a)=f.
\]

For a total endomorphism $e:Y\to Y$, put
$k=q_1e\Delta_X:X\to X$ and $N(e)=\langle q_0,kq_0\rangle:Y\to Y$.
Both morphisms are total. We have $q_0N(e)=q_0$, so $N(e)^2=N(e)$. If $e$ indexes
$f:Y\to Y$, precomposing its index equation with
$(\Delta_Xq_0)\times1_Y$ gives
$a^{-1}\tau(kq_0\times a)=fp_{1,Y,Y}$.
The left-hand side is $t(N(e)\times1_Y)$, so $N(e)$ is an
idempotent $t$-index of $f$.
\end{proof}

\section{Ranges and the diagonal domain}\label{sec:diagonal}

We first record the properties of ranges needed for the diagonal argument.

\subsection{Ranges}\label{sec:ranges}
In this subsection, ranges are defined in a restriction category with
restriction zeros and binary restriction coproducts. The definition
applies in particular to the $P$-categories under consideration.

For a morphism $\phi:X\to Y$ and $\varepsilon\in\mathrm{Dom}(Y)$, we write
$\phi\prec\varepsilon$ if and only if $\varepsilon\phi=\phi$; we
say that $\varepsilon$ {\it receives} $\phi$. Suppose that
$\varepsilon$ receives $\phi$ and that, for every object $Z$ and all
morphisms $\psi,\psi':Y\to Z$, the equality
$\psi\phi=\psi'\phi$ implies
$\psi\varepsilon=\psi'\varepsilon$. Then $\varepsilon$ is the least
domain in $\mathrm{Dom}(Y)$ receiving $\phi$. We call $\varepsilon$
the range of $\phi$ and write $\mathrm{ran}\phi=\varepsilon$.

The relation $\prec$ agrees with $\le$ on domains. It has the following transitivity property.

\begin{prop}\label{PropRec}
If $\MAP{\phi}{X}{Y}$ is a morphism and $\beta,\gamma$ are domains
in $Y$ such that $\phi\prec\beta$ and $\beta\le\gamma$, then
$\phi\prec\gamma$.
\end{prop}
\begin{proof}
We have
$\phi=\beta\phi=(\beta\gamma)\phi=\gamma(\beta\phi)=\gamma\phi$,
so that $\phi\prec\gamma$. 
\end{proof}

A restriction category has ranges if every morphism has a range in the sense above.
For a $P$-category $\mathbf C$, we require in addition that
for morphisms $\phi,\psi$ of $\mathbf{C}$,
$\mathrm{ran}(\phi\times\psi)=\mathrm{ran}\phi\times\mathrm{ran}\psi$,
so that ranges are compatible with the near product. Such a category has a calculus of ranges
~\cite{MR92e:03067,MR88i:03078}. 

In a restriction category with coproducts and
ranges, the meet-semilattice $\mathrm{Dom}(X)$ becomes a
distributive lattice if one defines the join by
$\varepsilon\vee\delta=\mathrm{ran}[\varepsilon,\ \delta]$ for
$\varepsilon,\delta\in\mathrm{Dom}(X)$
~\cite{MR88i:03078,MR92e:03067}. The next statement follows
from the definition of range ~\cite[Proposition
5.4]{MR88i:03078}.

\begin{prop}\label{PropDomRan}
If $\varepsilon$ is a domain in $\mathrm{Dom}(Y)$ and
$\MAP{\varphi}{X}{Y}$ is any morphism, then
$\mathrm{ran}(\varepsilon\varphi)=\varepsilon
\mathrm{ran}(\varphi)$.
\end{prop}
\begin{proof}
Put $r=\mathrm{ran}(\varphi)$ and $s=\mathrm{ran}(\varepsilon\varphi)$.
The domain $\varepsilon r$ receives $\varepsilon\varphi$, so
$s\le\varepsilon r$. Conversely, $s\varepsilon\varphi=\varepsilon\varphi$;
the equality property of $r$ gives $s\varepsilon r=\varepsilon r$.
Hence $\varepsilon r\le s$.
\end{proof}

The least-receiving property shows that
$\mathrm{ran}[\alpha,\beta]$ is the join of $\alpha$ and $\beta$.
The proposition gives distributivity, since
$\varepsilon(\alpha\vee\beta)=\mathrm{ran}[\varepsilon\alpha,
\varepsilon\beta]=\varepsilon\alpha\vee\varepsilon\beta$.

A domain $\gamma\in\mathrm{Dom}(X)$ is \textit{complemented} if
there is a domain $\gamma^\perp$ with
$\gamma\vee\gamma^\perp=1_X$ and $\gamma\gamma^\perp=0_X$.
Distributivity makes this complement unique.

\begin{prop}\label{PropExpected}
For any morphisms $\MAP{\varphi}{X}{Y},\MAP{\psi}{Y}{Z}$,
$\psi\varphi = 0$ if and only if
$\rst{\psi}\mathrm{ran}(\varphi)=0$.
\end{prop}
\begin{proof}
If $\rst{\psi}\mathrm{ran}(\varphi)=0$, then
$\psi\varphi=\psi\rst{\psi}\mathrm{ran}(\varphi)\varphi=0$.
Conversely, if $\psi\varphi=0$, then
$0=\rst{\psi\varphi}=\rst{\rst{\psi}\varphi}$ by the restriction identity $\rst{gf}=\rst{\rst g f}$, which implies that $\rst{\psi}\varphi=0$.
Take ranges and apply the preceding proposition.
\end{proof}

For a morphism $f:X\to Y$, define its \textit{inverse image}
$f^*:\mathrm{Dom}(Y)\to\mathrm{Dom}(X)$ and \textit{image}
$f_*:\mathrm{Dom}(X)\to\mathrm{Dom}(Y)$ by
$f^*\beta=\rst{\beta f}$ and $f_*\alpha=\mathrm{ran}(f\alpha)$.
Inverse image preserves meets, and image preserves joins. In a
$P$-category with coproducts and ranges, inverse image also preserves
binary joins \cite{MR88i:03078,MR92e:03067}. Whenever the target is
a restriction category without a specified near product, we require
this last property explicitly.
The following fact will be used below ~\cite[Cor. 5.6 (iii)]{MR88i:03078}.
\begin{prop}\label{PropImageFact}
Let $\MAP{\phi}{X}{Y}$ be a morphism, and let $\alpha\in\mathrm{Dom}(X)$ and $\beta\in\mathrm{Dom}(Y)$. Then
$\phi_*(\alpha)\beta =\phi_*(\alpha\phi^*\beta)$.
\end{prop}
\begin{proof}
By R4 and commutativity of domains,
$\phi\alpha\phi^*\beta=\beta\phi\alpha$.
Take ranges and apply Proposition~\ref{PropDomRan}.
\end{proof}

\begin{prop}
If $\MAP{\phi}{X}{Y}$ is total in a restriction category with ranges, then for each
$\alpha\in\mathrm{Dom}(X)$ and for each $\beta\in\mathrm{Dom}(Y)$,
$\phi_*\alpha\le\beta$ if and only if $\alpha\le\phi^*\beta$ \cite{MR92e:03067}.
\end{prop}

Adjointness follows from the next proposition.

\begin{prop}\label{propsemiadjointness}
If $\MAP{\phi}{X}{Y}$ is a morphism in a restriction category $\mathbf{C}$ with ranges, then for each
$\alpha\in\mathrm{Dom}(X)$ and for each $\beta\in\mathrm{Dom}(Y)$,

a) $\phi_*\alpha\le\beta$ implies  $\alpha\phi^*\beta=\alpha\rst{\phi}$; and

b) $\alpha\le\phi^*\beta$ implies $\phi_*\alpha\le\beta$.
\end{prop}

\begin{proof}
For a), the condition implies $\beta\mathrm{ran}(\phi\alpha)=\mathrm{ran}(\phi\alpha)$. Composing on the right with $\phi\alpha$, we have $\beta\phi\alpha=\phi\alpha$. Taking restrictions gives
$\rst{\beta\phi}\alpha=\rst{\beta\phi\alpha}
=\rst{\phi\alpha}=\rst\phi\alpha$.

Assertion b) follows from Proposition \ref{PropImageFact}. 
\end{proof}

\subsection{A domain with no complement}
\Needspace{12\baselineskip}
\begin{thm}\label{ThmNoncomp}
Let $\mathbf C$ be a $P$-recursion category with binary coproducts
and ranges. If $t:X\times X\to X$ is weak Turing and
$1_X\ne0_X$, then $\rst{t\Delta_X}$ is not complemented.
\end{thm}
\begin{proof}
Put $h=t\Delta_X$ and $\kappa=\rst h$. Suppose that a domain
$\varepsilon$ complements $\kappa$, and choose a total $t$-index
$e:X\to X$ of $\varepsilon$. Equation~\eqref{EqDiagonalIndex}
gives $he=\varepsilon e$, and hence
$e^*\kappa=\rst{he}=\rst{\varepsilon e}=e^*\varepsilon$.
Since $e$ is total and inverse image preserves meets and binary joins,
these equal domains have meet $0_X$ and join $1_X$. Therefore
$1_X=0_X$, a contradiction.
\end{proof}

A domain $\delta\in\mathrm{Dom}(X)$ is a \textit{quasicomplement}
of $\varepsilon\in\mathrm{Dom}(X)$ if $\delta\varepsilon=0_X$ and,
for every $\phi:Y\to X$, equality $\varepsilon\phi=\delta\phi$
implies $\phi=0$.
Di Paola and Heller proved that in a dominical recursion category the
diagonal domain has no quasicomplement
\cite[Corollary~8.6]{MR88i:03078}. Complements and quasicomplements
coincide in a dominical category.

\begin{prop}
In a $P$-category with a system of zeros, binary coproducts and ranges, the following assertions are equivalent.

(i) Weakly total morphisms are total.

(ii) ~\cite[Proposition 5.3]{MR88i:03078} If $\varepsilon\in\mathrm{Dom}(X)$ satisfies $\varepsilon\ne1_X$, there exists $\delta\in\mathrm{Dom}(X)$ with $\delta\ne0$ and $\varepsilon\delta=0$.

(iii) ~\cite[Proposition 7.3]{MR88i:03078} Complements and quasicomplements coincide.
\end{prop}

Rosolini has written in  ~\cite[page 313]{MR90e:18001} that in view of his representation Theorem 2.8 ({\it loc. cit.}), all the results of section 5 of ~\cite{MR88i:03078} ``...can be generalized to $P$-categories with ranges stable under products by checking that they hold in a $P$-category of the form
$\mathcal{M}$--$\mathrm{Ptl}(\mathbf{A})$.'' The ranges in $S$ are defined
by existential quantification of graph formulas \cite[p.~111]{MR90i:03046}
and are stable under products, since the two coordinate variables can
be quantified independently. The preceding proposition,
together with Montagna's weakly total non-total morphism in $S$
\cite[Example~1, p.~112]{MR90i:03046}, shows that Proposition~5.3 of
\cite{MR88i:03078} does not extend to this generality. Theorem~\ref{ThmNoncomp} gives the noncomplementation result directly.

\section{Recursively inseparable domains}\label{sec:inseparability}

Let $\mathbf C$ be a $P$-recursion category with coproducts and ranges.
A disjoint pair $(\alpha,\beta)$ in $\mathrm{Dom}(X)$ is
\textit{recursively separable} if there is a complemented domain
$\gamma\in\mathrm{Dom}(X)$ with $\alpha\le\gamma$ and
$\beta\le\gamma^\perp$. It is \textit{recursively inseparable}
if no such $\gamma$ exists.

\subsection{Connected domains}
Connected domains were introduced in \cite{MR92e:03067} and used
in \cite{MR93b:03103,leng}.
A domain $\alpha$ is \textit{connected} if every decomposition
$\alpha=\beta\vee\gamma$ with $\beta\gamma=0$ has $\beta=0$ or
$\gamma=0$. The zero domain is connected. A restriction category with
ranges and binary restriction coproducts is \textit{locally connected}
if, for all $f,g:X\to Y$, equality $f\alpha=g\alpha$ for every
connected domain $\alpha$ on $X$ implies $f=g$ \cite{MR92e:03067}.
When inverse images preserve binary joins, total morphisms preserve
connected domains under direct image. This gives the following
alternatives for a morphism into a coproduct.

\begin{prop}\label{PropAlternative}
Let $\mathbf C$ be a restriction category with ranges, restriction zeros,
and binary restriction coproducts, in which inverse images of domains
preserve binary joins. Suppose that $\alpha\in\mathrm{Dom}(X)$ is
connected.

\noindent(i) If $\MAP{f}{X}{Y}$ is total, then $f_*\alpha$ is connected.

\noindent(ii) If $\MAP{f}{X}{Y + Z}$ is total,
then $f_*\alpha\le 1_Y+0_Z$ or $f_*\alpha\le 0_Y+ 1_Z$.

\noindent(iii) If $\MAP{f}{X}{Y+ Z}$ is total,
then $f\alpha\prec 1_Y+0_Z$ or $f\alpha\prec 0_Y+1_Z$.
\end{prop}

\begin{proof}
The proof of (i) uses adjointness and Proposition \ref{PropImageFact}.
Suppose that $f_*\alpha=\beta\vee\gamma$ with $\beta\gamma=0$. Since $f$ is total, by
adjointness, $\alpha\le
f^*(\beta\vee\gamma)=f^*\beta\vee f^*\gamma$;
equivalently, $\alpha=\alpha f^*\beta\vee\alpha f^*\gamma$. We
have that
$\alpha(f^*\beta)\alpha(f^*\gamma)=\alpha f^*(\beta\gamma)=\alpha\rst{0f}=0$.
Therefore, either $\alpha f^*\beta=0$ or $\alpha f^*\gamma=0$.
If $\alpha f^*\beta=0$, then
$0=f_*(\alpha f^*\beta)=f_*(\alpha)\beta=(\beta\vee\gamma)\beta=\beta$.
Similarly, if $\alpha f^*\gamma=0$, $\gamma=0$.

For (ii), note that $f_*\alpha=f_*\alpha\wedge(1_Y+
1_Z)=f_*\alpha\wedge((1_Y+ 0)\vee(0+ 1_Z))
=(f_*\alpha\wedge(1_Y+ 0))\vee (f_*\alpha\wedge(0+ 1_Z))$;
the conclusion follows from (i).

For (iii), note that $f_*\alpha$ is a range, so that $f\alpha\prec
f_*\alpha$; i.e., $(f_*\alpha)f\alpha =f\alpha$. The result
follows from (ii) and Proposition \ref{PropRec}.
\end{proof}

\subsection{Range functors}
Let $\mathbf C$ and $\mathbf D$ be restriction categories with ranges,
restriction zeros, and binary restriction coproducts. A \textit{range
functor} $F:\mathbf C\to\mathbf D$ preserves restrictions, ranges,
zeros, and binary coproducts. It preserves total morphisms, meets and
joins of domains, and complements. For objects $X,Y$ of $\mathbf C$,
the canonical coproduct comparison is the isomorphism
\[
 \chi_{X,Y}=[Fi_0,Fi_1]:FX+FY\longrightarrow F(X+Y).
\]
For $P$-categories,
preservation of the near product is not required; examples occur in
categories of coalgebras \cite{leng}.

The \textit{kernel} of $F$ is
$\ker F=\{f:X\to Y\mid Ff=0_{FX,FY}\}$; it is closed under
composition on either side.

\begin{prop}\label{PropKernel}
If $F:\mathbf{C}\rightarrow\mathbf{D}$ is a range functor and if $\mathbf{D}$
is locally connected, then for each object $X$ in $\mathbf{C}$ such that $1_X\notin\ker F$, there exists a nonzero connected domain in $\mathrm{Dom}(FX)$.
\end{prop}

\begin{proof}
Suppose that $0$ is the only connected domain in $\mathrm{Dom}(FX)$. Then $1_{F X}0 = 0$, which implies that $F1_X = 1_{F X} = 0$, since $\mathbf{D}$ is locally connected. But then
$1_X\in\ker F$, contradiction.  
\end{proof}

\subsection{Sections and weak choice}
In a restriction category, a {\it section} of a morphism
$\MAP{\phi}{X}{Y}$  is a morphism $\MAP{\sigma}{Y}{X}$ such that
$\phi\sigma=\rst\sigma$ and $\phi\sigma\phi=\phi$. It follows that
$\rst{\sigma}=\mathrm{ran}(\phi)$. A
restriction category satisfies the {\it axiom of choice} if every morphism
has a section. A morphism $\phi:X\to Y$ is a \textit{partial monomorphism}
if, for every object $W$ and all morphisms $\theta,\theta':W\to X$,
the equality $\phi\theta=\phi\theta'$ implies
$\rst\phi\,\theta=\rst\phi\,\theta'$.
A restriction category satisfies the \textit{weak axiom of choice},
or \textit{weak choice}, if every partial monomorphism has a section.

To check that a partial recursive function $\phi$ is a partial monomorphism, it suffices to check that for all (total) recursive functions $\theta,\theta'$, $\phi\theta=\phi\theta'$ implies $\rst\phi\theta=\rst\phi\theta'$. This follows from the fact that a recursively enumerable set is either empty or the image of a total recursive function. The following is a generalization.

\begin{prop}\label{PropPartialMono}
Let $\mathbf{C}$ be a split restriction category, and let $\MAP{f}{X}{Y}$ be a morphism in $\mathbf{C}$. Suppose that, for every object $W$ of $\mathbf C$ and all total
morphisms $g,h:W\to X$, the equality $fg=fh$ implies
$\rst f\,g=\rst f\,h$. 
Then $f$ is a partial monomorphism.
\end{prop}

\begin{proof}
Let $g,h:W\to X$ satisfy $fg=fh$. Split the common restriction
$\varepsilon=\rst{fg}=\rst{fh}$ as $\varepsilon=me$, where
$e:W\to Z$, $m:Z\to W$, and $em=1_Z$. Since $m$ is a total monomorphism,
$\rst{fgm}=\rst{\varepsilon m}=\rst m=1_Z$.
Consequently $fgm$ and $fhm$ are total. 
Therefore
$g m$ and $h m$ are total. By the hypothesis, $\rst{f}g m = \rst{f}h m$. Moreover,
\[
\rst{f}g = \rst{f}g\rst{\rst{f}g} = \rst{f}g\rst{fg}
         = \rst{f}gme =\rst{f}hme =\rst{f}h\rst{fh}
         =\rst{f}h\rst{\rst{f}h} = \rst{f}h.
\]
Therefore, $f$ is a partial monomorphism.
\end{proof}

\begin{prop}\label{PropTotalCopMono}
Let $\mathbf C$ be a restriction category with ranges, restriction
zeros and binary restriction coproducts, in which inverse images of
domains along every morphism preserve binary joins. Let $X$ be an
object of $\mathbf C$, and suppose that $\alpha,\beta\in\mathrm{Dom}(X)$ satisfy $\alpha\beta=0_X$.
For every object $Y$ and all total morphisms $\theta,\theta':Y\to X+X$, equality
$\nabla_X(\alpha+\beta)\theta=\nabla_X(\alpha+\beta)\theta'$
implies $(\alpha+\beta)\theta=(\alpha+\beta)\theta'$.
\end{prop}
\begin{proof}
Put $p=\nabla_X(\alpha+\beta):X+X\to X$, so that
$\rst p=\alpha+\beta$. Let $i_0,i_1:X\to X+X$ be the coproduct
injections, and define
\[
 \begin{gathered}
 \rho_0=[1_X,0_X],\qquad \rho_1=[0_X,1_X]:X+X\to X,\\
 \varepsilon_i=i_i\rho_i:X+X\to X+X\qquad(i=0,1).
 \end{gathered}
\]
The summand domains $\varepsilon_0,\varepsilon_1$ are disjoint
and have join $1_{X+X}$. Define domains on $Y$ by
\[
 e_i=\theta^*\varepsilon_i,\qquad
 f_j=(\theta')^*\varepsilon_j,\qquad
 \delta_{ij}=e_i f_j\qquad(i,j\in\{0,1\}).
\]
Totality and inverse-image preservation give
$e_0\vee e_1=1_Y=f_0\vee f_1$. Distributivity gives
$\delta_{i0}\vee\delta_{i1}=e_i$, and hence the four domains
$\delta_{ij}$ have join $1_Y$.

For fixed $i,j\in\{0,1\}$, put
\[
 x_{ij}=\rho_i\theta\delta_{ij}:Y\to X,
 \qquad y_{ij}=\rho_j\theta'\delta_{ij}:Y\to X.
\]
The restriction identity
$\varepsilon_i\theta=\theta\rst{\varepsilon_i\theta}$ and its
analogue for $\theta'$ give
\[
 \theta\delta_{ij}=i_i x_{ij},\qquad
 \theta'\delta_{ij}=i_j y_{ij}.
\]
Write $a_0=\alpha$ and $a_1=\beta$. The hypothesis $p\theta=p\theta'$
gives the common morphism
\[
 z_{ij}=a_i x_{ij}=a_j y_{ij}:Y\to X.
\]
It follows that
\[
 (\alpha+\beta)\theta\delta_{ij}=i_i z_{ij},\qquad
 (\alpha+\beta)\theta'\delta_{ij}=i_j z_{ij}.
\]
If $i=j$, these morphisms agree. If $i\ne j$, then
$z_{ij}=a_i z_{ij}=a_i a_j y_{ij}=0_{Y,X}$, so both morphisms
are zero. Thus they agree for every pair $i,j$.

Form the copairing
\[
 q=[[[\delta_{00},\delta_{01}],\delta_{10}],\delta_{11}]
       :((Y+Y)+Y)+Y\longrightarrow Y.
\]
A domain receives $q$ precisely when it receives each $\delta_{ij}$,
so $\mathrm{ran}(q)=1_Y$. The preceding equalities give
$(\alpha+\beta)\theta q=(\alpha+\beta)\theta' q$.
The defining equality property of $\mathrm{ran}(q)=1_Y$ therefore gives
$(\alpha+\beta)\theta=(\alpha+\beta)\theta'$.
\end{proof}

\begin{prop}\label{PropCopOut}
Let $\mathbf C$ be a split restriction category with ranges,
restriction zeros and binary restriction coproducts. Suppose that inverse
images of domains along every morphism preserve binary joins and that
weak choice holds.
For every object $X$ and all $\alpha,\beta\in\mathrm{Dom}(X)$ with $\alpha\beta=0_X$,
the morphism $p=\nabla_X(\alpha+\beta):X+X\to X$ has a unique section
$f:X\to X+X$. Moreover, $pf=\rst f=\mathrm{ran}(p)=\alpha\vee\beta$.
\end{prop}
\begin{proof}
Propositions~\ref{PropTotalCopMono} and~\ref{PropPartialMono} show
that $p$ is a partial monomorphism; here $\rst p=\alpha+\beta$.
Weak choice supplies a section $f$. By R4 and the section equations,
$\rst p\,f=f\rst{pf}=f\rst f=f$.
If $f'$ is another section, then $pf=pf'=\mathrm{ran}(p)$;
partial monicity yields $f=\rst p\,f=\rst p\,f'=f'$.
Finally, $\mathrm{ran}(p)=\alpha\vee\beta$ by the definition of join.
\end{proof}

\begin{prop}\label{PropDichotomy}
Let $\mathbf C$ satisfy the hypotheses of Proposition~\ref{PropCopOut},
and let $\alpha\in\mathrm{Dom}(X)$ be complemented. The section
$f:X\to X+X$ of $p=\nabla_X(\alpha^\perp+\alpha)$ is total and satisfies $f\alpha=i_1\alpha$ and
$f\alpha^\perp=i_0\alpha^\perp$.
\end{prop}
\begin{proof}
Proposition~\ref{PropCopOut} gives $\rst f=1_X$ and $pf=1_X$.
Since $p(f\alpha)=\alpha=p(i_1\alpha)$, partial monicity gives
$\rst p\,f\alpha=\rst p\,i_1\alpha$.
The left-hand side is $f\alpha$, and the right-hand side is
$i_1\alpha$. The equation for $\alpha^\perp$ follows in the same way.
\end{proof}

\subsection{The inseparability theorem}
Di Paola and Heller proved the existence of
recursively inseparable domains in a dominical recursion category
$\mathbf{C}$ with coproducts, under the assumption that $\mathbf{C}$ satisfies the
axiom of choice ~\cite[Theorem 8.15]{MR88i:03078}. Rosolini removed the requirement that $\mathbf{C}$ be dominical and that products distribute over coproducts, 
provided domains split in $\mathbf{C}$ ~\cite{MR90e:18001}.

\begin{thm}\label{ThmRecIns}
Let $\mathbf C$ be a $P$-recursion category with binary coproducts and
ranges. Let $\mathbf D$ be a split restriction category with ranges,
restriction zeros and binary restriction coproducts, satisfying weak
choice, in which inverse images of domains along every morphism preserve
binary joins. Let $F:\mathbf C\to\mathbf D$ be a range functor,
$t:X\times X\to X$ a weak Turing morphism in $\mathbf C$ with
$1_X\notin\ker F$, and $b:X\to X+X$ an isomorphism.

For each complemented $\gamma\in\mathrm{Dom}(X)$, let
$s_\gamma:FX\to FX+FX$ be the unique section of
$\nabla_{FX}(F\gamma^\perp+F\gamma)$. Assume the \textit{lifting
hypothesis}: for every such $\gamma$, there exists a morphism
$d_\gamma:X\to X+X$ in $\mathbf C$ such that
$Fd_\gamma=\chi_{X,X}s_\gamma$. Then the domains
\[
\alpha=(bt\Delta_X)^*(1_X+0_X),\qquad
\beta=(bt\Delta_X)^*(0_X+1_X)
\]
form a recursively inseparable pair in $\mathrm{Dom}(X)$.
\end{thm}
\begin{proof}
Put $h=bt\Delta_X:X\to X+X$. Since inverse image preserves meets
and zero, $\alpha\beta=h^*((1_X+0_X)(0_X+1_X))=0_X$.
Suppose that a complemented domain $\gamma$ separates the pair. Put
$\Gamma=F\gamma$, $\Gamma'=F\gamma^\perp$, and $s=s_\gamma$.
The domains $\Gamma,\Gamma'$ on $FX$ are complementary, and
$s:FX\to FX+FX$ is total. Write $j_0,j_1:FX\to FX+FX$ for the target
coproduct injections. Proposition~\ref{PropDichotomy} gives
$s\Gamma=j_1\Gamma$ and $s\Gamma'=j_0\Gamma'$.

Let $P_0=1_{FX}+0_{FX}$ and $P_1=0_{FX}+1_{FX}$ be the summand
domains on $FX+FX$. The morphisms $P_0s$ and $j_0\Gamma'$ agree
after both $\Gamma$ and $\Gamma'$. Since
$\mathrm{ran}[\Gamma,\Gamma']=1_{FX}$, the defining equality property of this range gives
$P_0s=j_0\Gamma'$. Likewise $P_1s=j_1\Gamma$. Taking restrictions
gives
\[
 s^*P_0=\Gamma',\qquad s^*P_1=\Gamma.
\]

Choose $d_\gamma:X\to X+X$ with
$Fd_\gamma=\chi_{X,X}s$, and put
$u=b^{-1}d_\gamma:X\to X$.
Choose a total $t$-index $e:X\to X$ of $u$.
Equation~\eqref{EqDiagonalIndex} gives
$t\Delta_Xe=ue$, and hence $he=d_\gamma e$.
Put
\[
 H=\chi_{X,X}^{-1}Fh:FX\to FX+FX,
 \qquad E=Fe:FX\to FX.
\]
Applying $F$ to the preceding identity gives $HE=sE$.
The canonical comparison satisfies
$\chi_{X,X}P_0=F(1_X+0_X)\chi_{X,X}$ and
$\chi_{X,X}P_1=F(0_X+1_X)\chi_{X,X}$.
Since $F$ preserves restrictions, it follows that
$F\alpha=H^*P_0$ and $F\beta=H^*P_1$. Thus
\[
 \begin{aligned}
 E^*F\alpha&=(HE)^*P_0=(sE)^*P_0=E^*\Gamma',\\
 E^*F\beta&=(HE)^*P_1=(sE)^*P_1=E^*\Gamma.
 \end{aligned}
\]
The inequalities $F\alpha\le\Gamma$ and $F\beta\le\Gamma'$
therefore give
\[
 E^*\Gamma'\le E^*\Gamma,\qquad E^*\Gamma\le E^*\Gamma'.
\]
These two domains are equal, and their meet is
$E^*(\Gamma\Gamma')=0_{FX}$. Consequently both are zero.
The morphism $E$ is total, so the inverse-image join law gives
\[
 1_{FX}=E^*1_{FX}
       =E^*(\Gamma\vee\Gamma')
       =E^*\Gamma\vee E^*\Gamma'=0_{FX},
\]
contrary to $1_X\notin\ker F$.
\end{proof}

The lifting hypothesis holds if
$F:\mathbf C(X,X+X)\to\mathbf D(FX,F(X+X))$ is surjective. It also holds
when every complemented $\gamma\in\mathrm{Dom}(X)$ has a total
morphism $d_\gamma:X\to X+X$ with
$d_\gamma\gamma=i_1\gamma$ and
$d_\gamma\gamma^\perp=i_0\gamma^\perp$. Indeed,
$p_\gamma d_\gamma=1_X$, where
$p_\gamma=\nabla_X(\gamma^\perp+\gamma)$: both sides agree after
composition with $\gamma$ and $\gamma^\perp$, and
$\mathrm{ran}[\gamma,\gamma^\perp]=1_X$. Thus
$\chi_{X,X}^{-1}Fd_\gamma$ is a section of
$\nabla_{FX}(F\gamma^\perp+F\gamma)$, and uniqueness gives
$Fd_\gamma=\chi_{X,X}s_\gamma$.

\begin{prop}\label{PropDisjointJoins}
Let $\mathbf C$ be a restriction category with restriction zeros and
binary restriction coproducts. Suppose that every pair of parallel
morphisms with disjoint domains has a join. For each complemented
domain $\gamma:X\to X$, the morphism
$d_\gamma=(i_1\gamma)\vee(i_0\gamma^\perp):X\to X+X$
is total and satisfies $d_\gamma\gamma=i_1\gamma$ and
$d_\gamma\gamma^\perp=i_0\gamma^\perp$.
If $F:\mathbf C\to\mathbf D$ is a range functor and $\mathbf D$
satisfies the target hypotheses of Theorem~\ref{ThmRecIns}, then
$Fd_\gamma=\chi_{X,X}s_\gamma$, where $s_\gamma:FX\to FX+FX$
is the unique section of $\nabla_{FX}(F\gamma^\perp+F\gamma)$.
Thus $\chi_{X,X}s_\gamma$ has a preimage under $F$, as required
by the lifting hypothesis of
Theorem~\ref{ThmRecIns}.
\end{prop}
\begin{proof}
Since the coproduct injections are total,
$\rst{i_1\gamma}=\gamma$ and
$\rst{i_0\gamma^\perp}=\gamma^\perp$. The two morphisms therefore
have disjoint domains. Their upper-bound equations give
$d_\gamma\gamma=i_1\gamma$ and
$d_\gamma\gamma^\perp=i_0\gamma^\perp$, and their domains satisfy
$\gamma,\gamma^\perp\le\rst{d_\gamma}$. Hence
$1_X=\gamma\vee\gamma^\perp\le\rst{d_\gamma}$, so $d_\gamma$ is
total. The section argument above gives
$Fd_\gamma=\chi_{X,X}s_\gamma$.
\end{proof}

\subsection{Failure of the lifting hypothesis}\label{sec:counterexample}
A target section $s_\gamma:FX\to FX+FX$ can have no source
morphism $d_\gamma:X\to X+X$ with
$Fd_\gamma=\chi_{X,X}s_\gamma$, even when $F$ is a faithful range
functor that preserves the near product and reflects zero morphisms. We give a counterexample
in which the pair in Theorem~\ref{ThmRecIns} is separable.

Fix an acceptable enumeration $(\varphi_e)$ of the partial recursive
functions, and put $K=\{e:\varphi_e(e)\downarrow\}$. Let $(W_a^K)$
be the enumeration of domains of an acceptable numbering of partial
$K$-recursive functions. Fix a primitive recursive bijective pairing
function with projections $\pi_0,\pi_1$.

The objects of $\mathbf C_K$ are the sets generated from $\mathbb N$
by binary products and tagged coproducts, with their recursive codings.
Its morphisms $f:X\rightharpoonup Y$ are partial functions with
$K$-computably enumerable domains that admit ordinary partial recursive
extensions. Equality is equality of graphs, and composition is composition
of partial functions. We write $1_B$ for the partial identity on a
subset $B$ and use $\simeq$ for equality of partial values, including
simultaneous undefinedness. On $\mathbb N$, every morphism has the form
\[
f_{e,a}(x)\simeq
\begin{cases}
\varphi_e(x),&x\in W_a^K,\\
\uparrow,&x\notin W_a^K.
\end{cases}
\]
The total morphisms are exactly the ordinary
total recursive functions.

\Needspace{8\baselineskip}
\begin{prop}\label{PropCounterexample}
The category $\mathbf C_K$ is a $P$-recursion category with binary
coproducts and ranges. Its inclusion $F:\mathbf C_K\to\mathbf{Pfn}$
satisfies all hypotheses of Theorem~\ref{ThmRecIns} except the lifting
hypothesis; its diagonal pair can be recursively separable.
Moreover, $\mathbf C_K$ contains two parallel morphisms with disjoint
domains and no common upper bound in the restriction order.
\end{prop}
\begin{proof}
Write $D_f$ for the domain of a partial function $f$, and let $f_0,g_0$
be partial recursive extensions of composable morphisms $f,g$.
The composite has extension $g_0f_0$ and domain
$D_f\cap f_0^{-1}(D_g)$, which is $K$-computably enumerable.
Identities belong to the category, and associativity holds for partial
functions. Products, restrictions, and tagged copairings have the same
closure property. Thus the set-theoretic near product and coproduct give
the stated categorical structures. All objects are recursively isomorphic,
and the empty partial functions give the prodominical zeros.
The image of a $K$-computably enumerable set under a partial recursive
function is $K$-computably enumerable. Its partial identity is the range,
with the required equality property, and ranges commute with products.

Define $t(\langle e,a\rangle,x)=f_{e,a}(x)$. This morphism has a
partial recursive extension and a $K$-computably enumerable domain.
For $v:\mathbb N\times\mathbb N\rightharpoonup\mathbb N$, choose a
partial recursive extension $V(z,x)$, and write $D_v$ for the
$K$-computably enumerable domain of $v$. The ordinary and oracle
parameterization theorems \cite{MR90d:03072} give
ordinary total recursive functions $e(z),a(z)$ such that
$\varphi_{e(z)}(x)\simeq V(z,x)$ and
$W_{a(z)}^K=\{x:(z,x)\in D_v\}$.
Oracle parameterization substitutes $z$ into program code and is
ordinary recursive; the oracle is used by the indexed program.
Consequently $c(z)=\langle e(z),a(z)\rangle$ is total in $\mathbf C_K$
and $v=t(c\times1)$. Recursive object codings handle the other
parameter objects, proving the Turing property.

The inclusion into $\mathbf{Pfn}$ preserves restrictions, ranges,
coproducts, zeros, and near products. The target is split and locally
connected, its inverse images preserve joins, and injective partial
functions have partial inverses, so weak choice holds.

Take $b(2n)=i_0(n)$ and $b(2n+1)=i_1(n)$, and put
$p(q)\simeq\varphi_{\pi_0(q)}(q)$. The diagonal domains are
\[
\begin{aligned}
D_\alpha&=\{q:q\in W_{\pi_1(q)}^K,\ p(q)\downarrow
                   \text{ with an even value}\},\\
D_\beta&=\{q:q\in W_{\pi_1(q)}^K,\ p(q)\downarrow
                   \text{ with an odd value}\}.
\end{aligned}
\]
The set $A=\{q:p(q)\downarrow\text{ with an even value}\}$ is
computably enumerable, hence decidable relative to $K$. Both $A$ and
$A^c$ are $K$-computably enumerable, so their partial identities are
complementary domains of $\mathbf C_K$. Since
$D_\alpha\subseteq A$ and $D_\beta\subseteq A^c$, the pair is separable.

To identify the failed lift, fix an index $a_0$ and choose $e(n)$
effectively for a program that ignores its input, simulates
$\varphi_n(n)$, and returns $0$ on convergence. Then
$n\in K$ if and only if $\langle e(n),a_0\rangle\in A$, so $A$ is
nonrecursive. The morphism
$p_A=\nabla(1_{A^c}+1_A):\mathbb N+\mathbb N\rightharpoonup\mathbb N$
has a unique section in $\mathbf{Pfn}$, given by
\[
s_A(n)=
\begin{cases}
i_1(n),&n\in A,\\
i_0(n),&n\notin A.
\end{cases}
\]
This function and $b^{-1}s_A$ are total and nonrecursive. They therefore
have no preimages under the inclusion $F$.

The morphisms $f=i_1 1_A$ and $g=i_0 1_{A^c}$ from $\mathbb N$ to
$\mathbb N+\mathbb N$ belong to $\mathbf C_K$ and have disjoint
domains. If $d$ were a common upper bound, then $d1_A=f$ and
$d1_{A^c}=g$. Since $A\cup A^c=\mathbb N$, these equations force
$d=s_A$, contradicting the fact that every total morphism of
$\mathbf C_K$ is recursive. Thus the joins required by
Proposition~\ref{PropDisjointJoins} fail to exist, although the domains
of $\mathbf C_K$ have binary joins.
\end{proof}

\bigskip
\noindent\textbf{Acknowledgements.}
For their comments I thank Sergei Artemov, Robin Cockett, Alex Heller,
Alphonse Vasquez, and Noson Yanofsky. Robin Cockett suggested considering
union restriction categories in the inseparability argument.

\begin{thebibliography}{DPM91}
\bibitem[CL02]{MR1871071}
J.~R.~B. Cockett and Stephen Lack, \textit{Restriction categories. I.
Categories of partial maps}, Theoret. Comput. Sci. \textbf{270} (2002),
no.~1--2, 223--259.
\bibitem[CL03]{MR1963657}
J.~R.~B. Cockett and Stephen Lack, \textit{Restriction categories. II.
Partial map classification}, Theoret. Comput. Sci. \textbf{294} (2003),
no.~1--2, 61--102.
\bibitem[DPH87]{MR88i:03078}
Robert A. Di Paola and Alex Heller, \textit{Dominical categories: recursion
theory without elements}, J. Symbolic Logic \textbf{52} (1987), no.~3, 594--635.
\bibitem[DPM91]{MR93b:03103}
Robert A. Di Paola and Franco Montagna, \textit{Some properties of the
syntactic $P$-recursion categories generated by consistent, recursively
enumerable extensions of Peano arithmetic}, J. Symbolic Logic
\textbf{56} (1991), no.~2, 643--660.
\bibitem[Hel90]{MR92e:03067}
Alex Heller, \textit{An existence theorem for recursion categories},
J. Symbolic Logic \textbf{55} (1990), no.~3, 1252--1268.
\bibitem[Len04]{leng}
Florian Lengyel, \textit{More existence theorems for recursion categories},
Ann. Pure Appl. Logic \textbf{125} (2004), no.~1--3, 1--41.
\bibitem[Len26]{leng2026composition}
Florian Lengyel, \textit{Pointwise provable equality and the failure of
composition}, 2026, \href{https://arxiv.org/abs/2609.25556v2}{arXiv:2609.25556v2}.
\bibitem[Mon89]{MR90i:03046}
Franco Montagna, \textit{``Pathologies'' in two syntactic categories of
partial maps}, Notre Dame J. Formal Logic \textbf{30} (1989), no.~1, 105--116.
\bibitem[Odi89]{MR90d:03072}
Piergiorgio Odifreddi, \textit{Classical recursion theory}, North-Holland,
Amsterdam, 1989.
\bibitem[Ros86]{coneit}
G. Rosolini, \textit{Continuity and effectiveness in topoi}, D.Phil.
thesis, University of Oxford, 1986.
\bibitem[Ros88]{MR90e:18001}
G. Rosolini, \textit{Representation theorems for $P$-categories},
Categorical Algebra and its Applications (F. Borceux, ed.), Lecture Notes
in Mathematics, vol.~1348, Springer, Berlin, 1988, pp.~307--315.
\bibitem[Ste93]{MR95c:03110}
Stefano Stefani, \textit{A relativization mechanism in recursion categories},
J. Symbolic Logic \textbf{58} (1993), no.~4, 1251--1267.
\bibitem[Ste02]{MR2003a:18003}
S.~T. Stefani, \textit{An algebraic approach to categories of partial
morphisms}, J. Symbolic Logic \textbf{67} (2002), no.~1, 117--129.
\end{thebibliography}
\end{document}
